\chapter{Introducción}\label{cap1}

El presente documento expone el diseño detallado de un biorreactor tipo airlift,
elaborado bajo la metodología de diseño mecatrónico con el fin de crear el entorno
de experimentación óptimo para el estudio y caracterización del comportamiento de
un cultivo de hongos ascomicetos coprófilos.

Primeramente, se presentan los objetivos a cumplir con el desarrollo del prototipo,
seguido de la justificación, en la que se expone la problemática existente y el
alcance de este. El marco teórico por su parte presenta conceptos y términos que
conforman las bases teóricas para la comprensión del proyecto. En el estado del
arte se manejan algunos de los últimos avances de este tipo de tecnología que
mantienen relación con los objetivos planteados.

En la sección planteamiento del problema se desglosa a profundidad el problema
a resolver, además de describirse el sistema por medio de diagramas IDEF0, donde
se describe paso a paso las acciones a realizar para la obtención de una respuesta
favorable.

Por su parte el apartado diseño detallado se encuentra dividido en diversas
secciones, comenzando por la creación de la geometría general del tanque,
seguido del diseño del sistema de control de nutrientes, nivel y temperatura;
después la programación, etapa en la que se consideran los diagramas de flujo que
rigen el ciclo lógico del sistema, así como diagramas de estados y bosquejos de la
interfaz. La sección de diseño de electrónica recopila todos los elementos que
conforman las etapas de potencia, protocolos de comunicación y microcontrolador,
todo eso con su respectivo diseño de PCB. En selección de materiales se presentan
todos los elementos prefabricados que componen el sistema, así como una breve
explicación del porqué de su elección. Finalmente, el apartado plan de manufactura
recaba todos los procesos que se llevaran a cabo para la creación del prototipo.

\section{Objetivos}

\subsection{Objetivo General}

Diseñar y construir un biorreactor tipo airlift semiautomatizado que permita al
usuario realizar pruebas en hongos ascomicetos coprófilos, mediante el ajuste
previo de las variables temperatura, presión, nivel, y cantidad de nutrientes, así
como el monitoreo de oxigenación e iluminación, facilitando la investigación y
determinación de las condiciones óptimas para su cultivo.

\subsection{Objetivos específicos}

\begin{itemize}
  \item Diseñar y construir una estructura con dimensiones máximas de 80~cm $\times$ 80~cm $\times$ 100~cm, que cuente con compartimentos para la adición de los reactivos implicados en el crecimiento de los ascomicetos.
  \item Diseñar e implementar un algoritmo de control que gestione las transiciones entre etapas de crecimiento del hongo, ajustando la salida de los actuadores en función de las variables de entrada: temperatura, presión, oxigenación y cantidad de nutrientes.
  \item Diseñar e integrar un sistema electrónico que incorpore sensores, circuitos de acondicionamiento de señal y actuadores.
  \item Desarrollar e implementar una interfaz humano-máquina que muestre las variables críticas: temperatura, presión, oxigenación, cantidad de nutrientes, iluminación y nivel. Estas variables deberán actualizarse cada minuto y almacenarse en una base de datos descargable.
\end{itemize}
